%% 22장 최소제곱
\chapter{최소제곱}
\label{ch:least-squares}

\abstract{관측이 미지수보다 많은 연립방정식 $Ax\approx b$를 잔차의 제곱합이 가장 작도록 푸는 최소제곱은 회귀분석의 계산적 핵심이다. 이 장에서는 최소제곱 문제의 해를 기하적으로 특징짓는 정규방정식에서 출발하지만, 그것을 그대로 푸는 것이 왜 위험한지를 조건수로 설명한다. 정규방정식은 문제의 조건수를 제곱해 버리는 반면, 21장의 QR 분해로 풀면 그런 손실이 없다. 최소제곱 문제의 섭동 이론으로 잔차가 클 때 조건수의 제곱이 나타나는 것이 문제 자체의 성질임도 함께 확인한다. 이어서 열이 거의 일차종속인 랭크 부족 문제를 특이값 분해와 허용오차로 다루는 방법, 잔차를 더 높은 정밀도로 계산해 해를 고치는 반복 개선, 긴 행렬을 빠르게 분해하는 Cholesky QR을 다룬다. stardust의 회귀 영역이 이 도구들을 어떻게 조합하는지도 본다.}

\section{문제와 정규방정식}
\label{sec:normal-equations}

$A$가 $m\times n$이고 $m>n$이면 $Ax=b$는 대개 해가 없다. \term{최소제곱 문제}{least squares problem}는 잔차 $r=b-Ax$의 2-노름을 최소로 하는 $x$를 찾는다. 기하적으로 $Ax$는 $A$의 열 공간 위를 움직이므로, 최소는 $Ax$가 $b$의 그 공간 위로의 정사영일 때 얻어지고, 이때 잔차는 열 공간에 수직이다. 곧 $A^{\mathsf T}r=0$, 다시 말해
\begin{equation}
\label{eq:normal}
A^{\mathsf T}Ax=A^{\mathsf T}b
\end{equation}
이다. 이것이 \term{정규방정식}{normal equations}이다. $A$의 열이 일차독립이면 $A^{\mathsf T}A$는 대칭 양정치이므로 20장의 Cholesky 분해로 풀 수 있고, 연산도 $A^{\mathsf T}A$를 만드는 $mn^2$번이 대부분이어서 싸다. 그래서 오랫동안 교과서의 표준 해법이었다.

문제는 정확도이다. $A^{\mathsf T}A$의 2-노름 조건수는 $\kappa(A)^2$이다. 9장의 원리에 따라 후진 안정한 Cholesky로 풀어도 오차는 대략 $\kappa(A)^2u$이다. 반면 문제 자체의 조건수는, 뒤에서 보듯 잔차가 작을 때 대략 $\kappa(A)$이다. 정규방정식은 문제가 요구하지도 않는 손실을 만드는 것이다. 9장에서 본 분산의 교과서 공식이 데이터를 제곱하는 순간 조건수를 제곱했던 것과 같은 현상이다.

\section{QR로 풀기}
\label{sec:ls-qr}

$A=QR$을 21장의 Householder QR로 구하고 $Q$를 $m\times m$ 직교 행렬로 늘려 $Q=(Q_1,Q_2)$로 나누자. $Q_1$은 $A$의 열 공간의 정규직교 기저이고 $Q_2$는 그 직교 여공간의 기저이다. 2-노름은 직교 변환에서 변하지 않으므로
\[
\|b-Ax\|_2^2=\|Q^{\mathsf T}b-Q^{\mathsf T}Ax\|_2^2=\|Q_1^{\mathsf T}b-Rx\|_2^2+\|Q_2^{\mathsf T}b\|_2^2
\]
이다. 둘째 항은 $x$와 무관하므로 최소는 $Rx=Q_1^{\mathsf T}b$일 때이고, 그때의 잔차 노름이 $\|Q_2^{\mathsf T}b\|_2$이다. 삼각 방정식 하나로 해가 나오고, $A^{\mathsf T}A$는 만들지 않는다. $Q$를 명시적으로 만들 필요도 없다. Householder 반사들을 $b$에 차례로 적용하면 $Q^{\mathsf T}b$가 나오기 때문이다. 21장의 후진 오차 정리에 따라 이 해법은 후진 안정하다. 계산한 해는 $A$와 $b$를 각각 상대적으로 $u$ 정도 바꾼 문제의 정확한 최소제곱 해이다.

두 해법의 차이를 재어 보자. 조건수 $\kappa$인 $200\times10$ 행렬과 잔차가 $0$인 우변을 무작위로 만들어 풀면, 해의 상대오차가 \cref{tab:ls}와 같다. 정규방정식의 오차는 $\kappa^2u$를 따라 $\kappa=10^8$에서 해를 거의 다 잃지만, QR의 오차는 $\kappa u$ 정도이다.

\begin{table}[t]
\begin{sdtable}{colspec={ccc}, row{1}={font=\small\bfseries}}
\toprule
$\kappa(A)$ & 정규방정식(Cholesky) & Householder QR \\
\midrule
$10^4$ & $1.5\times10^{-9}$ & $8.5\times10^{-15}$ \\
$10^6$ & $1.7\times10^{-5}$ & $1.5\times10^{-12}$ \\
$10^8$ & $3.0\times10^{-1}$ & $1.3\times10^{-10}$ \\
\bottomrule
\end{sdtable}
\caption{잔차가 $0$인 $200\times10$ 최소제곱 문제에서 두 해법의 해의 상대오차. numpy로 binary64에서 계산했다.}
\label{tab:ls}
\end{table}

\section{섭동 이론}
\label{sec:ls-perturbation}

그렇다면 최소제곱 문제의 진짜 조건수는 얼마인가. $A$와 $b$를 상대적으로 $\epsilon$ 이하만큼 바꿀 때 해의 변화를 일차까지 따지면, $\theta$를 $b$와 $Ax$ 사이의 각이라 할 때
\[
\frac{\|\Delta x\|_2}{\|x\|_2}\le\epsilon\Bigl(\frac{2\kappa(A)}{\cos\theta}+\kappa(A)^2\tan\theta\Bigr)+O(\epsilon^2)
\]
이다\cite{wedin1972,golub2013matrix}. 잔차가 $0$이면 $\theta=0$이고 조건수는 $\kappa(A)$ 정도이다. 잔차가 크면 둘째 항의 $\kappa(A)^2$이 나타난다. 이것은 알고리즘의 탓이 아니라 문제 자체의 성질이다. 잔차가 클 때 $A$를 조금 바꾸면 열 공간이 조금 기울어지고, 큰 잔차 벡터의 그 기울어진 공간 위로의 사영이 해를 크게 움직이기 때문이다. 그래서 잔차가 큰 문제에서는 QR로 풀어도 $\kappa^2u$의 오차를 피할 수 없지만, 잔차가 작은 문제에서 정규방정식만이 쓸데없이 $\kappa^2u$를 만든다. 9장의 원리 ``전진 오차 $\lesssim$ 조건수 $\times$ 후진 오차''가 여기서도 그대로 적용된다.

\section{랭크 부족}
\label{sec:rank-deficient}

열이 일차종속이거나 거의 일차종속이면 최소제곱 해가 하나로 정해지지 않거나, 정해지더라도 아주 크고 불안정하다. 이때는 해 가운데 노름이 가장 작은 것을 고르는 것이 표준이다. 24장에서 볼 특이값 분해 $A=U\Sigma V^{\mathsf T}$를 쓰면 \term{최소 노름 해}{minimum-norm solution}는
\[
x=\sum_{i:\,\sigma_i>\tau}\frac{u_i^{\mathsf T}b}{\sigma_i}v_i
\]
이다. 여기서 문턱 $\tau$보다 작은 특이값은 $0$으로 본다. 정확한 산술에서는 $\tau=0$이지만, 반올림이 있으면 수학적으로 $0$인 특이값이 $u\sigma_1$ 정도의 작은 값으로 계산되므로 4장에서 본 수치적 랭크의 허용오차가 필요하다. 표준적인 선택은 $\tau=\max(m,n)\,u\,\sigma_1$ 정도이다. 문턱이 너무 작으면 반올림이 만든 가짜 특이값이 $1/\sigma_i$로 크게 증폭되어 해를 오염시키고, 너무 크면 실제 정보를 버린다. 열 피벗 QR(21장)로 랭크를 판정하는 더 싼 방법도 있지만, 21장의 Kahan 행렬이 보여 주듯 확실하지는 않다.

\begin{instardust}[\code{docs/design/linalg.md} §10.4]
stardust의 \code{lstsq}는 특이값 분해 위의 분할 정복 해법인 LAPACK의 \code{gelsd}를 쓴다. 허용오차 \code{rtol}은 ``가장 큰 특이값의 \code{rtol}배 이하인 특이값을 $0$으로 본다''로 정의하고 기본값은 $\max(m,n)$배의 기계 엡실론이다. 이 정의를 끝까지 지키는 데서 몇 가지 미묘한 문제가 드러났다. LAPACK은 \code{rcond}가 $1$ 이상이면 조용히 엡실론으로 바꾸므로, stardust는 \code{rtol}$\ge1$이면 정의대로 랭크 $0$, 해 $0$을 직접 낸다. 또 \code{gelsd}는 크기가 $25$를 넘으면 분할 정복 전에 아주 작은 원소를 엡실론 크기로 바꾸어 정확히 $0$인 특이값마저 $0$이 아니게 만드므로, \code{rtol}이 엡실론보다 작으면 랭크가 \code{matrix\_rank}와 같도록 특이값 분해로 직접 푼다. 설계 문서에는 numpy와 비트까지 같으려면 우변의 메모리 주소가 $16$바이트 경계에서 같은 자리에 있어야 했다는, Accelerate의 \code{dgelsd}에 대한 관찰까지 기록되어 있다.
\end{instardust}

\section{반복 개선}
\label{sec:ls-refinement}

QR의 해는 $\kappa u$ 정도의 오차를 지닌다. 이것을 더 줄이려면 해를 고치는 단계를 덧붙인다. 계산한 해 $\hat x$의 잔차 $\hat r=b-A\hat x$를 더 높은 정밀도로 계산하고, 그 잔차에 대한 보정을 같은 분해로 풀어 더한다. 연립방정식에서는 단순히 $A\,\delta=\hat r$을 풀면 되지만, 최소제곱 문제에서는 잔차와 해가 함께 움직이므로 Björck가 제안한 대로 둘을 함께 미지수로 둔 \term{증강 시스템}{augmented system}
\[
\begin{pmatrix}I&A\\A^{\mathsf T}&0\end{pmatrix}\begin{pmatrix}r\\x\end{pmatrix}=\begin{pmatrix}b\\0\end{pmatrix}
\]
의 잔차를 계산해 보정한다\cite{bjorck1967iterative}. 증강 시스템의 보정도 QR 분해 하나로 풀 수 있다. 잔차를 10장의 double-double로 계산하면 보정 몇 번 만에 해가 binary64로 반올림한 정확한 해에 가까워진다.

\begin{instardust}[\code{docs/design/regression.md} §3.1]
stardust의 선형 회귀(\code{ols}, \code{wls}, \code{gls})는 가중치로 백색화한 설계 행렬을 QR로 분해해 첫 해를 구한 뒤, 잔차를 double-double로 계산해 같은 QR로 증강 시스템 꼴의 보정을 푼다. 보정은 다음 단계에서 의미가 없어질 때까지 많아야 세 번 한다. 설계 문서의 정확도 약속은 ``설계 행렬의 조건수가 $10^{11}$ 이하이면 계수 벡터의 상대오차가 $10^{-14}$ 이하''이다. 열이 겹치면, 곧 피벗 없는 QR에서 $|r_{jj}|$가 그 열의 노름의 \code{rtol}배 이하이면 그 열을 빼고 다시 분해하는데, 앞 열을 남기는 이 선택은 R의 회귀 함수와 같다. 빠진 열의 계수와 표준오차는 NaN으로 돌려준다.
\end{instardust}

\section{Cholesky QR}
\label{sec:cholesky-qr}

정규방정식의 속도와 QR의 정확도를 함께 얻는 방법도 있다. 행이 열보다 훨씬 많은 긴 행렬에서 Householder QR은 반사를 하나씩 적용하는 2층 BLAS 연산이 대부분이어서 느리다. \term{Cholesky QR}{Cholesky QR}은 $A^{\mathsf T}A=R^{\mathsf T}R$을 Cholesky로 분해하고 $Q=AR^{-1}$로 두는데, 둘 다 3층 연산이라 빠르지만 $Q$의 직교성이 $\kappa(A)^2u$만큼 무너진다. 그런데 이렇게 얻은 $Q$는 이미 거의 직교이므로 조건수가 작고, 같은 일을 $Q$에 한 번 더 하면 직교성이 $u$ 수준으로 회복된다. 이것을 \term{CholeskyQR2}{CholeskyQR2}라 한다\cite{fukaya2014cholqr2}. 첫 단계의 Cholesky가 실패하지 않으려면 $\kappa(A)^2u$가 충분히 작아야 하므로, 조건수가 대략 $u^{-1/2}\approx10^8$보다 작을 때만 쓸 수 있다\cite{yamamoto2015}.

\begin{instardust}[\code{docs/design/regression.md} §3.1]
stardust의 선형 회귀는 쓰인 행 $m$과 열 $k$가 $m\ge2k$, $mk\ge16384$일 때 먼저 CholeskyQR2를 시도한다. 첫 인수 $R_1$에서 어림한 조건수가 $10^5$ 이하이고 겹치는 열의 판정에 걸리는 열이 없을 때만 그 결과를 쓰고, 아니면 Householder QR로 돌아간다. 이론적인 한계 $10^8$보다 훨씬 보수적인 문턱인데, 설계 문서는 그 이유를 Arb 참값과의 비교로 적어 둔다. 문턱 안에서 표준오차와 지렛대의 오차가 Householder 판과 같거나 더 작았다. 긴 행렬에서 QR이 약 $4$배, 적합 전체가 약 $2$배 빨라졌다고 한다.
\end{instardust}

\begin{aside}[Longley의 경고]
1967년 Longley는 미국 노동통계국의 거시경제 자료로 회귀를 하면서 당시의 통계 프로그램들을 비교했다. 여섯 개의 설명 변수가 강하게 상관된 이 자료에서, 많은 프로그램이 정규방정식을 단정밀도로 풀어 유효숫자 하나도 맞지 않는 계수를 내놓았다\cite{longley1967}. 이 논문은 통계 소프트웨어의 수치 정확도에 대한 관심을 일깨웠고, Longley 자료는 지금도 미국 국립표준기술연구소의 회귀 참값 자료집에 포함되어 회귀 프로그램을 시험하는 데 쓰인다\cite{nist-strd}.
\end{aside}

\section{맺음}

최소제곱은 9장의 원리를 가장 선명하게 보여 주는 문제이다. 정규방정식은 대수적으로는 올바르지만 데이터를 제곱하면서 조건수까지 제곱하고, QR은 같은 문제를 문제의 조건이 허락하는 정확도로 푼다. 랭크가 부족하면 허용오차로 수치적 랭크를 정해야 하고, 더 높은 정확도가 필요하면 잔차를 정밀하게 계산해 고친다. 다음 장에서는 선형대수의 또 다른 큰 문제인 고유값 문제로 넘어간다.

최소제곱 문제의 수치적 방법은 Björck의 책\cite{bjorck1996}이 가장 포괄적이고, Lawson과 Hanson의 책\cite{lawson1995}은 고전이다.

\chapterreferences
